\documentclass[10pt,a4paper]{amsart}
\usepackage[margin=19mm]{geometry}
\usepackage{amsmath,amssymb,mathtools}
\usepackage[scr=rsfs,cal=euler]{mathalfa}
\usepackage{xfrac}
\usepackage[unicode,hidelinks,bookmarks=false]{hyperref}
\newcommand{\ie}{i.e.\ }
\newcommand{\cf}{cf.\ }
\newcommand{\resp}{respectively\ }
\newcommand{\Subsection}{\subsection}
\providecommand{\nobreakdash}{\nobreak\hbox{-}}
\setlength{\emergencystretch}{2em}
\multlinegap0pt
\usepackage{amssymb}
\usepackage{tikz}
\newtheorem{theorem}{Theorem}
\newtheorem{proposition}{Proposition}
\newtheorem{corollary}{Corollary}
\title{Higher-order Volterra-type integral operator on Hardy and Bergman spaces}
\author{Rahim Kargar}
\date{Source: arXiv:2512.16412v4 (CC BY 4.0)}
\begin{document}
\maketitle

\noindent\emph{Source and licence.} Higher-order Volterra-type integral operator on Hardy and Bergman spaces. Version arXiv:2512.16412v4 (CC BY 4.0), distributed by arXiv under the Creative Commons Attribution 4.0 International licence, http://creativecommons.org/licenses/by/4.0/, as recorded in the arXiv licence metadata for this exact version and shown on the abstract page https://arxiv.org/abs/2512.16412v4. Full licence notice: \texttt{licenses/arxiv-2512.16412-CC-BY-4.0.txt}.\par\smallskip

\noindent\textit{Editorial scope.} Counted unit: the article's corollary cor:not\_invariant\_n (source0.tex lines 1107--1126). The article's complete original proof of it is delivered with it (source0.tex lines 1128--1188, one \texttt{proof} environment of 61 lines). No mathematics is written, repaired or re-derived: the statement and the proof are the article's own lines, in the article's own notation, and no retained line is substituted or re-flowed.  Retained uncounted as the source-local context the proof needs: lines 934--951; lines 1072--1080.  Nothing was omitted from any retained span, so the package carries no omission record.  No retained line contains a citation, so the wrapper carries no bibliography and no bibliography entry.  No retained line contains a figure environment, an image-inclusion command or drawing code, so no image is delivered and the package declares no asset dependency.  Editorial formatting only, in addition: an AMS \texttt{amsart} preamble that restates the article's own declarations and macros used by the retained lines, namely the environments \texttt{theorem}, \texttt{proposition}, \texttt{corollary} on separate counters, together with the standard packages those lines need.  The retained environments print in this document as Theorem 1, Proposition 1, Corollary 1 in order of appearance, whereas the article prints the same items with the numbering of its own full document.  The complete native source of the article as distributed in the arXiv e-print for this exact version is bundled unchanged as \texttt{sources/191237/source0.tex}.
\begin{theorem}
\label{thm:bergman-boundedness}
Let $n \in \mathbb{N}$, $1 \le p \le q < \infty$, and $ \alpha, \beta > -1 $ satisfy the balance condition
\begin{equation}\label{eq:balance}
\frac{q}{p} \le \frac{\beta + 2}{\alpha + 2}.
\end{equation}
Let $ g $ be analytic in $ \mathbb{D} $.
Then the following are equivalent:
\begin{enumerate}
\item[(i)] $ T_{g,n}: A_\alpha^p \to A_\beta^q $ is bounded;
\item[(ii)] The measure
\begin{equation*}
d\mu_n(z) = |g'(z)|^q (1 - |z|^2)^{\beta + nq} \, dA(z)
\end{equation*}
is a $(p,q)$-Carleson measure for $ A_\alpha^p $.
\end{enumerate}
Moreover, $ T_{g,n}: A_\alpha^p \to A_\beta^q $ is compact if and only if $ \mu_n $ is a vanishing $(p,q)$-Carleson measure.
\end{theorem}

\begin{proposition}\label{prop:carleson-condition}
Let $1 \leq p \leq q < \infty$, $\alpha > -1$, and $r \in (0,1)$.
Let $\mu$ be a positive Borel measure on $\mathbb{D}$ that is a $(p,q)$-Carleson measure for $A_\alpha^p$.
Then, for every $a \in \mathbb{D}$,
\begin{equation*}
\mu(D(a,r)) \leq C (1 - |a|^2)^{\frac{q(2+\alpha)}{p}},
\end{equation*}
where $C>0$ depends only on $\alpha$, $p$, $q$, and $r$.
\end{proposition}

\begin{corollary}\label{cor:not_invariant_n}
Let $1\le p\le q<\infty$ and $\alpha,\beta>-1$ satisfy
\begin{equation*}
\frac{q}{p}\le \frac{\beta+2}{\alpha+2}.
\end{equation*}
Then the boundedness (and likewise compactness) of
\begin{equation*}
T_{g,n}:A_\alpha^p\to A_\beta^q
\end{equation*}
is, in general, not invariant under changes of the order $n$.

More precisely, there exist analytic symbols $g$ and integers $n\ge2$ such that
\begin{equation*}
T_{g,n}:A_\alpha^p\to A_\beta^q \ \text{is bounded (respectively compact),}
\end{equation*}
while
\begin{equation*}
T_{g,1}:A_\alpha^p\to A_\beta^q \ \text{fails to be bounded (respectively compact).}
\end{equation*}
\end{corollary}

\begin{proof}
By Theorem~\ref{thm:bergman-boundedness}, boundedness of $T_{g,n}$ is equivalent to
$\mu_n$ being a $(p,q)$-Carleson measure for $A_\alpha^p$, where
\begin{equation*}
d\mu_n(z)=|g'(z)|^q(1-|z|^2)^{\beta+nq}\,dA(z).
\end{equation*}
For $n\ge2$,
\begin{equation*}
d\mu_n(z)=(1-|z|^2)^{q(n-1)}\,d\mu_1(z)\le d\mu_1(z) \quad \text{since } 0<1-|z|^2<1.
\end{equation*}
Hence, if $\mu_1$ is a $(p,q)$-Carleson measure (respectively a vanishing $(p,q)$-Carleson
measure), then so is $\mu_n$. Therefore, boundedness (respectively compactness) of
$T_{g,1}$ implies boundedness (respectively compactness) of $T_{g,n}$ for every $n\ge2$.

To see that the converse may fail, fix $n\ge2$ and consider a symbol $g$ with
\begin{equation*}
g'(z)=(1-z)^{-\gamma},\quad \gamma>0.
\end{equation*}
Choose $\gamma$ so that
\begin{equation}\label{eq:gamma-interval}
\frac{\beta+q+2}{q}-\frac{2+\alpha}{p}
\ <\ \gamma \ \le\
\frac{\beta+nq+2}{q}-\frac{2+\alpha}{p}.
\end{equation}
Such a $\gamma$ exists because $n\ge2$ makes the right endpoint larger than the left endpoint by $n-1>0$.

Let $r\in(0,1)$ be fixed and let $D(a,r)$ denote the Bergman ball. For $a$ tending to $1$ and $z\in D(a,r)$ one has $|1-z|\asymp |1-a|\asymp 1-|a|\asymp 1-|a|^2$ on $D(a,r)$, hence
\begin{equation*}
|g'(z)|^q = |1-z|^{-\gamma q}\asymp (1-|a|^2)^{-\gamma q},
\quad z\in D(a,r),
\end{equation*}
and also $(1-|z|^2)\asymp (1-|a|^2)$ on $D(a,r)$. Therefore,
\begin{align*}
\mu_n(D(a,r))
&=\int_{D(a,r)} |g'(z)|^q(1-|z|^2)^{\beta+nq}\,dA(z)\\
&\asymp (1-|a|^2)^{-\gamma q}\,(1-|a|^2)^{\beta+nq}\,A(D(a,r))
\asymp (1-|a|^2)^{\beta+nq+2-\gamma q}.
\end{align*}
By Proposition \ref{prop:carleson-condition}, the $(p,q)$-Carleson condition for $A_\alpha^p$ is equivalent to the estimate
\begin{equation*}
\mu(D(a,r))\lesssim (1-|a|^2)^{q(2+\alpha)/p}\quad (a\in\mathbb D),
\end{equation*}
and, hence $\mu_n$ is a $(p,q)$-Carleson measure if and only if
\begin{equation*}
\beta+nq+2-\gamma q \ \ge\ \frac{q(2+\alpha)}{p}.
\end{equation*}
Similarly, $\mu_1$ is a $(p,q)$-Carleson measure if and only if
\begin{equation*}
\beta+q+2-\gamma q \ \ge\ \frac{q(2+\alpha)}{p}.
\end{equation*}
By the choice \eqref{eq:gamma-interval}, $\mu_n$ is $(p,q)$-Carleson but $\mu_1$ is not.
Therefore, $T_{g,n}$ is bounded while $T_{g,1}$ fails to be bounded (and hence is not
compact), which proves the stated non-invariance.

Finally, by choosing $\gamma$ so that
\begin{equation*}
\beta+nq+2-\gamma q > \frac{q(2+\alpha)}{p}
\end{equation*}
while still violating the corresponding condition for $n=1$, one obtains that $\mu_n$
is a vanishing $(p,q)$-Carleson measure, yielding the compactness statement.
\end{proof}

\end{document}
